\section{Operating Agents}
\sectionkicker{AGENT OPERATIONS}
\begin{wideflow}
{\Large\bfseries エージェント運用――ハーネス最適化とデプロイ・セキュリティボット}\par
\smallskip
{\color{SurveyMuted}9月23日のCursorの二つの動きを、測られた数値の出所と切り分けて示す。}\par
\end{wideflow}

9月23日、Cursorは利用者のトークン代を7\%ほど軽くしたと発表した\cite{cursoreff}。大勢の利用での試し（A/B）によるものとされ、質を落とさずにとの留保がつく。中身は、システムプロンプトを3分の2ほど削り、ツール定義を使るときだけ読む構成に替え（固定の説明トークンが6割減、MCPツールでは先行して46.9\%減）、プロンプトキャッシュの区切りを安定した層の後に置いてコールドキャッシュミスを2割減らし、読みの行番号を10行ごとに間引いて1.6\%ほど浮かせ、サブエージェント利用の促しを外して釣り合いを取った、とされる。いずれもCursorの実運用環境での測りであり、他のハーネスにそのまま移せる数ではない。proxy指標としての評価は当てにならず、実運用の利用分布を見る姿勢も添えられる。

同じ日、デプロイ工程の仕上げにRolloutsとSecurity Reviewerの二つのボットが出た\cite{cursorrollouts}。RolloutsはPull Request（PR）から本番までを見守り、回帰不具合を見つけては知らせ、止め、ロールバック案を置くとされる。今日のところ自らマージもロールバックもしない。Security ReviewerはPRごとに読み、攻撃経路と修正案を示すとされ、レビューの平均時間が4.8分から3.8分、受け入れが45～50\%から60～70\%に動いたと図に示される。いずれも作り手の測りであり、他所での再現はない。使えるのはTeams・Enterpriseプランで、10日間のトライアルが添えられる。機能フラグ（feature flag）の上げ下げやリリーストレインへの目配りはこれからとされる。

\begin{claimboundary}[資料と評価の境界]
7\%はCursorの実運用環境の数であり、他のハーネスへの一般化には使わない。見直しの効果の数値は作り手の報告であり、独立した再現を待つ。これからの構成は出荷済みとして扱わない。春先に公開されたセキュリティエージェントのテンプレートとは別の製品である。
\end{claimboundary}
